\chapter{FUTURAS LÍNEAS DE TRABAJO}

La implementación actual del pipeline se ha desplegado mediante Docker Compose, lo cual resulta adecuado para el alcance de este TFM (un entorno de desarrollo y validación en una única máquina), pero presenta limitaciones evidentes de cara a un despliegue productivo real: no ofrece escalado horizontal automático de los componentes con mayor carga (particiones de Kafka, ejecutores de Spark Structured Streaming).

Una línea de evolución natural del sistema es la migración de la orquestación hacia una distribución ligera de Kubernetes, manteniendo el resto de la arquitectura (Mosquitto, Kafka en modo KRaft, Apicurio y los sinks duales TimescaleDB/PostgreSQL) sin cambios sustanciales, ya que estos componentes ya están containerizados y son, en principio, portables a manifiestos de Kubernetes o a Helm charts. Spark Structured Streaming y Kafka requieren cambios antes de la migración: el primero debe empaquetarse como imagen Docker (hoy corre desde el entorno virtual del host), y el segundo debe ampliarse a varios \textit{brokers} para ofrecer tolerancia real a fallos, hoy limitada por su factor de replicación de uno.

Entre las distribuciones de orquestación evaluadas como candidatas destacan dos por su bajo consumo de recursos y su idoneidad para escenarios de IoT y edge computing:

\begin{itemize}
    \item \textbf{K3s} (Rancher/SUSE): una distribución certificada por la CNCF (Cloud Native Computing Foundation), lo que asegura compatibilidad completa con la API estándar de Kubernetes, que sustituye \textit{etcd} por SQLite en despliegues de un solo nodo y reduce drásticamente el consumo de memoria frente a Kubernetes estándar.
    \item \textbf{MicroK8s} (Canonical): distribución con fuerte integración en el ecosistema Ubuntu (el mismo utilizado en el entorno de desarrollo de este proyecto) y un sistema de complementos (\textit{addons}) que simplifica la activación de componentes como almacenamiento persistente o ingress.
\end{itemize}

Como trabajo futuro, se propone:

\begin{enumerate}
    \item Migrar el stack de Docker Compose a manifiestos de Kubernetes (o Helm charts), comenzando por un clúster de un solo nodo sobre K3s o MicroK8s en el mismo entorno WSL/Ubuntu ya utilizado para el desarrollo.
    \item Evaluar el escalado horizontal de los ejecutores de Spark Structured Streaming y de las particiones de Kafka bajo carga variable, midiendo consumo de CPU, memoria y latencia, replicando la metodología de H. Koziolek y N. Eskandani \cite{koziolek2023}.
    \item Explorar la incorporación de mecanismos de auto-recuperación y despliegue continuo (GitOps) sobre el clúster resultante, como extensión natural del repositorio actual del proyecto.
\end{enumerate}
